\section{Conclusiones y Trabajos Futuros}

\subsection{Conclusiones Principales}

Tras desarrollar el marco completo, queda clara una convicción: la convergencia nativa entre conectividad Direct-to-Device y redes no terrestres representa mucho más que una extensión de cobertura; constituye un cambio paradigmático necesario para materializar la visión de conectividad ubicua y resiliente del ecosistema 6G.

El trabajo ha demostrado que es posible diseñar una arquitectura híbrida TN-NTN que combine payloads regenerativos, orquestación inteligente y soporte D2D multi-banda, logrando mejoras significativas en disponibilidad de servicio y flexibilidad operativa. Al comparar con propuestas existentes de conectividad integrada, el marco presentado aporta una visión más holística que une estándares 3GPP, desafíos del canal y consideraciones operacionales de forma coherente \cite{Ullah2026_D2D}.

Lejos de pretender resolver todos los interrogantes del campo, este artículo establece una base técnica reproducible que equilibra ambición con realismo de implementación. Los análisis de trade-offs y las soluciones propuestas para Doppler, movilidad y gestión de recursos ofrecen un punto de partida sólido para operadores, fabricantes y organismos de estandarización.

\subsection{Direcciones de Investigación Futuras}

Uno de los desafíos persistentes al concluir un estudio de esta naturaleza consiste en separar las extensiones prometedoras de aquellas que, aunque interesantes, podrían diluir el esfuerzo principal.

Resulta particularmente llamativo el potencial de evolucionar hacia arquitecturas que integren de forma nativa comunicación y posicionamiento preciso en el mismo enlace D2D-NTN \cite{DeGaudenzi2026_NTN}. Esta línea permitiría habilitar nuevos servicios de localización de alta precisión en entornos donde los sistemas GNSS tradicionales fallan o resultan insuficientes.

Otra dirección prioritaria radica en profundizar la inteligencia de la orquestación TN-NTN mediante aprendizaje continuo y gemelos digitales distribuidos. Revisiones recientes del campo sugieren que solo mediante optimización conjunta y predictiva se podrá alcanzar la escalabilidad masiva y la eficiencia energética que 6G demanda \cite{Wang2025_NTN6G}.

Quedan además pendientes temas críticos como la validación a gran escala en entornos reales, el desarrollo de terminales de bajo consumo optimizados para uplink satelital, la estandarización de mecanismos de handover zero-touch y el análisis profundo de implicaciones regulatorias y de soberanía de datos en constelaciones globales.

En síntesis, aunque el camino hacia una convergencia D2D-NTN madura y comercialmente viable sigue siendo exigente, el marco técnico propuesto aquí ofrece una contribución concreta y un conjunto claro de caminos para avanzar. Creemos firmemente que esta integración se convertirá en uno de los habilitadores fundamentales de las redes 6G del futuro.